\section{Unequal centers: normal convergence through every order jet}
\label{clc:sec:unequal}

The preceding finite identities share one affine center $A$.
For independent centers $A_j$, define
\begin{equation}\label{clc:eq:U}
 \UU_{\bq,\bA}(\bs;\mu)=
 \int_{\R^{r-1}}\prod_{j=1}^r h_{q_j,A_j,s_j}(x_j)
             \dd\boldsymbol x.
\end{equation}
We require a common weight:
\begin{equation}\label{clc:eq:unequal-strip}
 \max\bigl(0,-\min_j\Rea A_j\bigr)<c<q_*.
\end{equation}
This condition is stronger than the separate requirements
$\Rea A_j>-q_j$ when the scales are unequal; it must not be dropped.
It holds, in particular, if all $\Rea A_j>-q_*$.
The weighted proof of Theorem~\ref{clc:thm:addition} applies to give
\begin{equation}\label{clc:eq:U-contour}
 \UU_{\bq,\bA}(\bs;\mu)=\frac1{2\pi i}
 \int_{c-i\infty}^{c+i\infty}e^{p\mu}G_{\bq}(p)
                \prod_{j=1}^r(p+A_j)^{-s_j}\dd p.
\end{equation}
It is entire in its independent orders; no integer-order restriction is
being analytically differentiated.

Choose a real reference $A$ with $A+c>0$ and set $\delta_j=A_j-A$.
Suppose
\begin{equation}\label{clc:eq:rho-condition}
 \rho_0=\max_j\frac{|\delta_j|}{A+c}<1.
\end{equation}
Define explicit polynomial coefficients by
\begin{align}
 \prod_{j=1}^r(1+\delta_jz)^{-s_j}
   &=\sum_{N=0}^\infty h_N(\bs;\bdelta)z^N,\label{clc:eq:hN-gen}\\
 h_N(\bs;\bdelta)
   &=\sum_{n_1+\cdots+n_r=N}
       \prod_{j=1}^r\frac{(-\delta_j)^{n_j}(s_j)_{n_j}}{n_j!}.
 \label{clc:eq:hN}
\end{align}
In particular $h_0=1$, $h_1=-\sum_js_j\delta_j$, and
\[
 h_2=\frac12\sum_js_j(s_j+1)\delta_j^2
                 +\sum_{i<j}s_is_j\delta_i\delta_j.
\]

\begin{theorem}[Normally convergent unequal-center closure]
\label{clc:thm:unequal}
Under \eqref{clc:eq:unequal-strip} and \eqref{clc:eq:rho-condition},
\begin{equation}\label{clc:eq:unequal-series}
 \boxed{\displaystyle
 \UU_{\bq,\bA}(\bs;\mu)
 =\sum_{N=0}^\infty h_N(\bs;\bdelta)
                  \CC_{\bq,A}(W+N;\mu),\qquad W=\sum_js_j.}
\end{equation}
The series is normally convergent on compact subsets where a strict
version of \eqref{clc:eq:rho-condition} holds, jointly in the independent
orders, the centers, and $\mu$ in a closed substrip of
$|\Ima\mu|<\pi\sum_jq_j^{-1}$. Every fixed mixed order or center
derivative may be taken term by term. Center collisions are regular.
If the scales are commensurate, every coefficient function $\CC$ on the
right has the finite pole-free Hurwitz or Lerch form already proved.
\end{theorem}
\begin{proof}
On the contour, $|p+A|\ge A+c$, so
$|\delta_j/(p+A)|\le\rho_0<1$ uniformly. The chosen logarithms satisfy
\[
 \log(p+A_j)=\log(p+A)+
             \log\left(1+\frac{\delta_j}{p+A}\right).
\]
This is first clear as $\delta_j$ moves from zero inside the strict disk,
and continues there because both $p+A$ and $p+A_j$ stay in the right
half-plane. The binomial series therefore gives
\[
 \prod_j(p+A_j)^{-s_j}
 =(p+A)^{-W}\sum_{N\ge0}h_N(\bs;\bdelta)(p+A)^{-N}.
\]
On a compact set of orders, $|(s)_n/n!|$ is at most a fixed polynomial
in $n$; this follows directly from the Gamma ratio away from a finite
initial set, with continuity covering nonpositive integer orders.
The absolute binomial series is consequently bounded by a convergent
polynomial-times-geometric series, uniformly along the whole contour.
The remaining factors in \eqref{clc:eq:U-contour} have an integrable
exponential vertical majorant, locally uniformly in all parameters.
Tonelli's theorem for the majorant and ordinary dominated convergence
justify termwise contour integration, giving \eqref{clc:eq:unequal-series}.

The same estimates hold on a slightly larger compact neighborhood of
any strict parameter set. Cauchy's integral formula then proves normal
convergence of every fixed derivative, not merely of the undifferentiated
series. In particular no singularity is created when two $\delta_j$
coincide. The final assertion follows by substituting
Theorem~\ref{clc:thm:finite} into each term.
\end{proof}

\begin{corollary}[Every common-strip configuration has such an expansion]
\label{clc:cor:global-center}
For any fixed centers satisfying \eqref{clc:eq:unequal-strip}, a reference
$A$ satisfying \eqref{clc:eq:rho-condition} exists. One may choose a real
number larger than zero and all the finitely many numbers
\begin{equation}\label{clc:eq:choose-A}
 \frac{|A_j|^2-c^2}{2(c+\Rea A_j)}.
\end{equation}
\end{corollary}
\begin{proof}
The denominators are positive. The strict inequality
$|A_j-A|^2<(A+c)^2$ is exactly
$|A_j|^2-c^2<2A(c+\Rea A_j)$, so the stated choice works for every $j$.
\end{proof}

This proves the unequal-shift research question in the preceding
reciprocal report \cite[Further research, question 2]{RLC} for its
right-half-plane shifts, and extends it to unequal scales whenever the
common strip exists. For equal scales $q_j=1$, the old shifts are
$a_j=1+A_j$; their conditions $\Rea a_j>0$ imply that such a common
$c<1$ can be chosen. Thus the original equal-scale question is fully
covered, including all independent spectral derivatives.

The result should not be described as a finite reduction of arbitrary
unequal centers: the outer series is generally infinite. The finite
closure is retained \emph{inside each term}. Theorem
\ref{clc:thm:alignment} below explains why generic unequal centers do
not recover pure order addition.
